\section*{The Appointed Propers in Full}

\noindent The Latin below is the complete recurring formulary between marginal
nos.~1562 and 1571 in the 1962 Vatican typical Missal, printed pp.~393--394.
It was visually checked against the identified CMAA facsimile; the surviving
pre-1962 text is reproduced as received, without modernization. Scriptural
English is the public-domain Douay--Rheims/Challoner; the Collect, Secret, and
Postcommunion are the public-domain English of the Cummiskey hand missal
(Philadelphia, 1861). Neither historical English witness is official 1962
liturgical text. Where a chant differs from the Clementine text translated by
the Douay, the gap is disclosed rather than filled by an editorial translation.

\fulltextelement{Introit\enspace\properrefs{Int.}}{\latin{Psalmus 73:20, 19 et 23; 73:1}; no.~1562}
Réspice, Dómine, in testaméntum tuum, et ánimas páuperum tuórum ne
derelínquas in finem: exsúrge, Dómine, et iúdica causam tuam, et ne
obliviscáris voces quæréntium te. \textnormal{Ps. ibid., 1} Ut quid, Deus,
repulísti in finem: irátus est furor tuus super oves páscuæ tuæ?
℣.~Glória Patri.
\finishfulltextelement
\begin{fulltextenglish}{Douay--Rheims/Challoner, Ps.~73:20, 19, 22--23 and 1}
Have regard to thy covenant: for they that are the obscure of the earth have
been filled with dwellings of iniquity. Deliver not up to beasts the souls
that confess to thee: and forget not to the end the souls of thy poor. Arise,
O God, judge thy own cause: remember thy reproaches with which the foolish man
hath reproached thee all the day. Forget not the voices of thy enemies: the
pride of them that hate thee ascendeth continually. O God, why hast thou cast
us off unto the end: why is thy wrath enkindled against the sheep of thy
pasture? ℣.~\latin{Glória Patri.}
\end{fulltextenglish}
\fulltextgap{The chant has \latin{Dómine} where the source verses do not,
\latin{derelínquas} where the Clementine has \latin{obliviscáris}, and
\latin{voces quæréntium te} where it has the voices of enemies. The printed
reference also omits v.~22 although the antiphon uses it. The complete Douay
verses expose those differences; no composite English antiphon is supplied.}

\fulltextelement{Collect\enspace\properrefs{Coll.}}{\latin{Oratio}; no.~1563}
Omnípotens sempitérne Deus, da nobis fidei, spei et caritátis augméntum: et,
ut mereámur ássequi quod promíttis, fac nos amáre quod prǽcipis. Per Dóminum.
\finishfulltextelement
\begin{fulltextenglish}{Cummiskey hand missal (1861), Collect, p.~423}
O Almighty and eternal God, grant us an increase of faith, hope, and charity;
and, that we may deserve what thou promisest, make us love what thou
commandest. Thro'.
\end{fulltextenglish}
\fulltextgap{The English does not separately render \latin{ássequi}, ``to
attain.'' The commentary therefore reasons from the Latin when it distinguishes
loving the command from attaining the promise.}

\fulltextelement{Epistle\enspace\properrefs{Ep.}}{\latin{Ad Galatas 3:16--22}; no.~1564}
Léctio Epístolæ beáti Pauli Apóstoli ad Gálatas.

Fratres: Abrahæ dictæ sunt promissiónes, et sémini eius. Non dicit: Et
semínibus, quasi in multis; sed quasi in uno: Et sémini tuo, qui est Christus.
Hoc autem dico: testaméntum confirmátum a Deo, quæ post quadringéntos et
trigínta annos facta est lex, non írritum facit ad evacuándam promissiónem.
Nam si ex lege heréditas, iam non ex promissióne. Abrahæ autem per
repromissiónem donávit Deus. Quid ígitur lex? Propter transgressiónes pósita
est, donec veníret semen, cui promíserat, ordináta per Angelos in manu
mediatóris. Mediátor autem uníus non est: Deus autem unus est. Lex ergo
advérsus promíssa Dei? Absit. Si enim data esset lex, quæ posset vivificáre,
vere ex lege esset iustítia. Sed conclúsit Scriptúra ómnia sub peccáto, ut
promíssio ex fide Iesu Christi darétur credéntibus.
\finishfulltextelement
\begin{fulltextenglish}{Douay--Rheims/Challoner, Gal.~3:16--22}
To Abraham were the promises made and to his seed. He saith not: And to his
seeds as of many. But as of one: And to thy seed, which is Christ. Now this I
say: that the testament which was confirmed by God, the law which was made
after four hundred and thirty years doth not disannul, to make the promise of
no effect. For if the inheritance be of the law, it is no more of promise.
But God gave it to Abraham by promise. Why then was the law? It was set
because of transgressions, until the seed should come to whom he made the
promise, being ordained by angels in the hand of a mediator. Now a mediator is
not of one: but God is one. Was the law then against the promises of God: God
forbid! For if there had been a law given which could give life, verily
justice should have been by the law. But the scripture hath concluded all
under sin, that the promise, by the faith of Jesus Christ, might be given to
them that believe.
\end{fulltextenglish}
\fulltextgap{The address \latin{Fratres} is a liturgical incipit borrowed from
v.~15 and is not part of the translated verses. The lesson stops at v.~22;
Paul's pedagogue image begins later and is context, not appointed text.}

\fulltextelement{Gradual\enspace\properrefs{Grad.}}{\latin{Psalmus 73:20, 19 et 22}; no.~1565}
Réspice, Dómine, in testaméntum tuum: et ánimas páuperum tuórum ne
obliviscáris in finem. ℣.~Exsúrge, Dómine, et iúdica causam tuam: memor esto
oppróbrii servórum tuórum.
\finishfulltextelement
\begin{fulltextenglish}{Douay--Rheims/Challoner, Ps.~73:20a, 19b and 22}
Have regard to thy covenant \dots{} and forget not to the end the souls of thy
poor. ℣.~Arise, O God, judge thy own cause: remember thy reproaches with which
the foolish man hath reproached thee all the day.
\end{fulltextenglish}
\fulltextgap{The closing \latin{memor esto oppróbrii servórum tuórum} is not
the Clementine's Ps.~73:22b; it corresponds verbally, apart from the vocative,
to Ps.~88:51. The Missal nevertheless cites Ps.~73. The Douay is not altered to
hide that boundary.}

\fulltextelement{Alleluia\enspace\properrefs{All.}}{\latin{Psalmus 89:1}; no.~1566}
Allelúia, allelúia. ℣.~Ps. 89, 1 Dómine, refúgium factus es nobis a generatióne et
progénie. Allelúia.
\finishfulltextelement
\begin{fulltextenglish}{Douay--Rheims/Challoner, Ps.~89:1}
Alleluia, alleluia. ℣.~Lord, thou hast been our refuge from generation to
generation. Alleluia.
\end{fulltextenglish}
\fulltextgap{The chant's \latin{a generatióne et progénie} is not the
Clementine's \latin{a generatióne in generatiónem}; the Douay translates the
latter.}

\fulltextelement{Gospel\enspace\properrefs{Gosp.}}{\latin{Lucas 17:11--19}; no.~1567}
✠ Sequéntia sancti Evangélii secúndum Lucam.

In illo témpore: Dum iret Iesus in Ierúsalem, transíbat per médiam Samaríam et
Galilǽam. Et cum ingrederétur quoddam castéllum, occurrérunt ei decem viri
leprósi, qui stetérunt a longe; et levavérunt vocem dicéntes: Iesu præcéptor,
miserére nostri. Quos ut vidit, dixit: Ite, osténdite vos sacerdótibus. Et
factum est, dum irent, mundáti sunt. Unus autem ex illis, ut vidit quia
mundátus est, regréssus est, cum magna voce magníficans Deum, et cécidit in
fáciem ante pedes eius, grátias agens: et hic erat Samaritánus. Respóndens
autem Iesus, dixit: Nonne decem mundáti sunt? et novem ubi sunt? Non est
invéntus qui redíret, et daret glóriam Deo, nisi hic alienígena. Et ait illi:
Surge, vade; quia fides tua te salvum fecit.
\finishfulltextelement
\begin{fulltextenglish}{Douay--Rheims/Challoner, Luke 17:11--19}
And it came to pass, as he was going to Jerusalem, he passed through the midst
of Samaria and Galilee. And as he entered into a certain town, there met him
ten men that were lepers, who stood afar off. And lifted up their voice,
saying: Jesus, Master, have mercy on us. Whom when he saw, he said: Go, shew
yourselves to the priests. And it came to pass, as they went, they were made
clean. And one of them, when he saw that he was made clean, went back, with a
loud voice glorifying God. And he fell on his face before his feet, giving
thanks. And this was a Samaritan. And Jesus answering, said: Were not ten made
clean? And where are the nine? There is no one found to return and give glory
to God, but this stranger. And he said to him: Arise, go thy way; for thy faith
hath made thee whole.
\end{fulltextenglish}
\par\smallskip\latin{Credo.}
\fulltextgap{The Missal supplies \latin{In illo témpore} and the subject
\latin{Iesus}; the Clementine verse and its Douay English carry the subject
from the preceding narrative. The Creed itself belongs to the Ordinary.}

\fulltextelement{Offertory\enspace\properrefs{Off.}}{\latin{Psalmus 30:15--16}; no.~1568}
In te sperávi, Dómine; dixi: Tu es Deus meus, in mánibus tuis témpora mea.
\finishfulltextelement
\begin{fulltextenglish}{Douay--Rheims/Challoner, Ps.~30:15--16}
But I have put my trust in thee, O Lord: I said: Thou art my God. My lots are
in thy hands. Deliver me out of the hands of my enemies; and from them that
persecute me.
\end{fulltextenglish}
\fulltextgap{The Missal sings \latin{témpora mea}, ``my times,'' where the
Clementine reads \latin{sortes meæ}, translated ``my lots.'' This Vulgate
Ps.~30:15--16 is Ps.~31:15--16 when the title is numbered as verse~1, as in
the NABRE, and Ps.~31:14--15 when the title is unnumbered. No editorial English
replaces the registered witness.}

\fulltextelement{Secret\enspace\properrefs{Sec.}}{\latin{Secreta; Præfatio de Sanctissima Trinitate}; no.~1569}
Propitiáre, Dómine, pópulo tuo, propitiáre munéribus: ut, hac oblatióne
placátus, et indulgéntiam nobis tríbuas, et postuláta concédas. Per Dóminum
nostrum.

Præfatio de Ssma Trinitate.
\finishfulltextelement
\begin{fulltextenglish}{Cummiskey hand missal (1861), Secret, p.~424}
Be thou propitious, O Lord, to thy people: and mercifully receive their
offerings: that being appeased thereby, thou mayest grant us pardon, and hear
our requests. Thro'.
\end{fulltextenglish}
\fulltextgap{The Latin repeats \latin{propitiáre}; the English renders the
second imperative by ``mercifully receive.'' The printed \latin{Ssma}
contracts \latin{Sanctissima}.}

\fulltextelement{Communion\enspace\properrefs{Comm.}}{\latin{Sapientia 16:20}; no.~1570}
Panem de cælo dedísti nobis, Dómine, habéntem omne delectaméntum, et omnem
sapórem suavitátis.
\finishfulltextelement
\begin{fulltextenglish}{Douay--Rheims/Challoner, Wis.~16:20}
Instead of which things, thou didst feed thy people with the food of angels,
and gavest them bread from heaven, prepared without labour; having in it all
that is delicious, and the sweetness of every taste.
\end{fulltextenglish}
\fulltextgap{The antiphon adapts rather than quotes the verse: it addresses
God in the first person plural, changes the verb, omits ``prepared without
labour,'' and recasts both closing phrases. The complete biblical verse is
therefore printed without manufacturing an English chant text.}

\fulltextelement{Postcommunion\enspace\properrefs{Postcomm.}}{\latin{Postcommunio}; no.~1571}
Sumptis, Dómine, cæléstibus sacraméntis: ad redemptiónis ætérnæ, quæsumus,
proficiámus augméntum. Per Dóminum.
\finishfulltextelement
\begin{fulltextenglish}{Cummiskey hand missal (1861), Postcommunion, p.~425}
May these heavenly mysteries, O Lord, which we have received, advance our
eternal redemption. Thro'.
\end{fulltextenglish}
\fulltextgap{In the Latin the communicants are the subject of
\latin{proficiámus}, and \latin{augméntum} repeats the Collect's noun. The
historical English changes the grammatical subject and leaves the repeated
noun implicit; the commentary follows the Latin.}
